\section{Translation: MIL as a coarsening problem}
\label{sec:translation}

\subsection{The DGP}

Consider a domain of instances partitioned into $K$ \emph{instance
types}; the continuous-feature extension, in which $K$ is replaced
by a feature map, is developed in \cref{sec:continuous}. Each type
$k$ has an unknown instance-positive probability
\begin{equation}
  \eta_k = \Prob\{y = 1 \mid \text{instance of type } k\}
  \in [0,1).
  \label{eq:trans-eta}
\end{equation}
A dataset consists of $N$ bags. Bag $i$ contains a multiset of
instances summarized by its \emph{composition vector}
$\bm{m}_i \in \mathbb{Z}_{\geq 0}^K$, where $m_{ik}$ is the count of
type-$k$ instances in bag $i$. Conditional on the composition, the
instance labels $y_{ij}$ are independent Bernoulli$(\eta_{k(j)})$,
and the observed bag label is the logical-OR
\begin{equation}
  Y_i = \max_j y_{ij}
  = \1\Big\{ \textstyle\sum_j y_{ij} \geq 1 \Big\}.
  \label{eq:trans-or}
\end{equation}
We observe $(\bm{m}_i, Y_i)$ together with instance features; the
instance labels $y_{ij}$ are latent. Writing the per-type
log-survival $s_k = -\log(1 - \eta_k) \geq 0$ and stacking
$\bm{s} = (s_1,\ldots,s_K)$, independence gives the closed-form bag
likelihood
\begin{equation}
  \Prob\{Y_i = 0 \mid \bm{m}_i\}
  = \prod_{k=1}^K (1 - \eta_k)^{m_{ik}}
  = \exp\!\big(-\bm{m}_i^\top \bm{s}\big).
  \label{eq:trans-lik}
\end{equation}
Equation \eqref{eq:trans-lik} says the negative-log-probability of a
negative bag is \emph{additive} in instance contributions: bag
labels follow a Bernoulli model with log-survival link
$g(\mu) = -\log(1-\mu)$ and
linear predictor $\bm{m}_i^\top \bm{s}$.

\subsection{The translation}

\Cref{tab:translation} states the explicit correspondence to
masked-data inference.

\begin{table}[ht]
\centering
\small
\begin{tabular}{ll}
\toprule
\textbf{Masked-data framework} & \textbf{Multiple instance learning} \\
\midrule
Component $j \in \{1,\ldots,m\}$ & Instance type $k \in \{1,\ldots,K\}$ \\
Component event (failure) & Instance label $y_{ij} = 1$ \\
System failure $T_i = \min$ / OR of events & Bag label $Y_i = \max_j y_{ij}$ \\
Failed cause $K_i$ & Which instance is positive (latent) \\
\textbf{Candidate set} $c_i$ & \textbf{Bag} (its instance composition $\bm{m}_i$) \\
Singleton candidate $|c_i| = 1$ & Singleton bag (an instance-level label) \\
Identifying functions & Instance feature representation \\
Masking mechanism & Bag-formation mechanism \\
C1 (cause in candidate set) & $Y_i = 1 \Rightarrow$ a positive instance is in bag $i$ \\
C2 (symmetric masking) & Bag assignment is exchangeable in instance labels \\
C3 (parameter independence) & Bag formation does not depend on $\bm{\eta}$ \\
\bottomrule
\end{tabular}
\caption{Translation between the masked-data framework
\citep{towell2026masked} and multiple instance learning.}
\label{tab:translation}
\end{table}

One structural feature has no counterpart in the earlier
applications. In scRNA-seq, the ambiguous observation was a
\emph{zero} count; in MIL, the ambiguous observation is a
\emph{positive} bag. A negative bag resolves every instance label
($y_{ij} = 0$ for all $j$, by C1) and so is fully informative,
whereas a positive bag is the coarsened observation. The
masked-data framework is polarity-agnostic, so the apparatus
transfers unchanged; the practical consequence is that the
identifying ``singletons'' here are positive instance-level labels,
not the known-rate controls used in scRNA-seq.

Under deterministic OR, C1 holds by construction: a positive bag
contains a positive instance with probability one. C2 holds when
the rule for grouping instances into bags does not look at the
latent labels, and C3 holds when that grouping does not depend on
$\bm{\eta}$. Both are reasonable for benchmark MIL, where bags are
fixed by the data-collection protocol (a molecule, an image, a
slide) rather than constructed adaptively.

\subsection{Existing methods in this language}

\paragraph{Diverse Density \citep{maron1998framework}} maximizes a
noisy-OR bag likelihood over a target concept point. In our
language, this is maximum likelihood under \eqref{eq:trans-lik}
with $\eta$ parametrized by distance to the concept point and a
firing rate below one (see \cref{sec:identifiability} for the
resulting confound).

\paragraph{mi-SVM and MI-SVM \citep{andrews2003support}} impute
instance labels subject to the bag-OR constraint and fit a
max-margin instance classifier. The OR constraint is exactly C1;
the imputation step is a hard-assignment surrogate for the
candidate-set sum.

\paragraph{Attention-based deep MIL \citep{ilse2018attention}} pools
instance embeddings with learned attention weights. The attention
weights are commonly read as instance importances; the rank
condition of \cref{sec:identifiability} states precisely when that
reading estimates a parameter.

\paragraph{CLAM \citep{lu2021data}} applies attention MIL to
whole-slide pathology images and adds instance-level clustering
supervision. The clustering supervision injects partial
instance-level labels: singleton candidate sets in our framework.

These methods are not in conflict; they are different
parametrizations and different surrogates for the same coarsening
problem. The framework lets us state when each is identifiable and
what bias it inherits when its aggregation assumption fails.
